\chapter{Verification of the Hand-Written Layers}
\label{ch:verification}

\section{Why verification is methodologically necessary}

This chapter is placed before the results, rather than in an appendix, because its content is a
precondition for those results meaning anything.

The argument is as follows. Consider a hand-written LSTM cell in which the input and forget gates
have been transposed --- a single-character error in the order of a tuple unpacking. The resulting
model is not broken in any observable sense. It compiles, it trains, its loss decreases, and it
reports a test accuracy perhaps two or three points below what a correct implementation would
achieve. Nothing in the output distinguishes it from a correct implementation of a slightly weaker
architecture.

Now consider the purpose of this report: to compare two architectures and to interpret the
resulting difference. A three-point degradation introduced by an implementation defect in one of
the two models is of precisely the magnitude that would be reported, in good faith, as an
architectural finding. The experiment would produce a confident, well-presented and entirely
spurious conclusion.

\keypoint{The characteristic failure mode of from-scratch numerical code is silent degradation,
not an observable error.} Testing is therefore not a matter of engineering hygiene here; it is
what licenses the inferences in Chapter~\ref{ch:results}.

Nine checks are specified below. All are executed on every run of the notebook, before any training
occurs, and the suite raises an assertion on failure rather than emitting a warning, so that it is
not possible to obtain results from an unverified implementation.

\section{The reference implementations}

Two library primitives are used in this chapter and nowhere else in the project:
\code{nn.LSTM} and \code{F.layer\us{}norm}. Their use here is legitimate, and indeed necessary: the
strongest available evidence that a hand-written layer is correct is that it agrees numerically
with an independent, widely-tested implementation of the same mathematics.

The distinction being drawn is between \emph{using} a primitive to build a model --- which the
from-scratch constraint forbids --- and \emph{comparing against} one to verify a model, which is
the only rigorous way to discharge that constraint.

\section{Checks 1--2: equivalence against reference implementations}

\subsection{Check 1 --- the LSTM cell}

This is the strongest check in the suite. Rather than testing that the hand-written cell behaves
\emph{similarly} to a reference, it tests that it is \emph{identical} given identical parameters.

The procedure copies the reference implementation's input-projection matrix, its hidden-projection
matrix, and the \emph{sum} of its two bias vectors into the hand-written cell, then runs both over
the same random input sequence and compares every timestep.

{\small
\begin{verbatim}
ref_lstm = nn.LSTM(D, H, num_layers=1, batch_first=True)   # reference only
ours = LSTMLayerScratch(D, H, reverse=False)
with torch.no_grad():
    ours.cell.W.weight.copy_(ref_lstm.weight_ih_l0)
    ours.cell.W.bias.copy_(ref_lstm.bias_ih_l0 + ref_lstm.bias_hh_l0)
    ours.cell.U.weight.copy_(ref_lstm.weight_hh_l0)
ref_out, (ref_h, _) = ref_lstm(x)
our_out, our_h = ours(x, full_mask)
\end{verbatim}}

Observed maximum absolute deviation: $5.96\times10^{-8}$ across all timestep outputs and
$4.47\times10^{-8}$ on the final hidden state. Both are at the level of single-precision rounding.

This result simultaneously confirms the gate algebra of
Equations~\ref{eq:lstm_i}--\ref{eq:lstm_h}, the gate \emph{ordering} convention, the handling of
the fused bias, and the correctness of the unrolled recurrence. A transposition of any two gates
would produce a deviation of order unity.

\subsection{Check 2 --- layer normalisation}

The hand-written normalisation is compared against \code{F.layer\us{}norm} with the same learned
parameters and a deliberately non-trivial input distribution (mean 1.5, standard deviation 3.0), so
that an error in either centring or scaling would be visible.

Observed maximum absolute deviation: $4.77\times10^{-7}$. This confirms in particular the use of
the biased variance estimator; substituting the unbiased estimator produces a deviation several
orders of magnitude larger.

\section{Check 3: attention against an independent formulation}

The head split and merge is the most error-prone part of a multi-head attention implementation,
because the reshaping operations are compact and their failure mode is a silent permutation of
features rather than an exception.

A naive reference is therefore computed, which processes one \emph{(batch, head)} pair at a time
and slices head columns explicitly:

{\small
\begin{verbatim}
for b in range(x.size(0)):
    head_outputs = []
    for h in range(n_heads):
        sl = slice(h * d_k, (h + 1) * d_k)
        qh, kh, vh = q_all[b, :, sl], k_all[b, :, sl], v_all[b, :, sl]
        s = (qh @ kh.T) / math.sqrt(d_k)
        s = s.masked_fill(~km[b].unsqueeze(0), -1e9)
        head_outputs.append(torch.softmax(s, dim=-1) @ vh)
    naive[b] = mha.W_o(torch.cat(head_outputs, dim=-1))
\end{verbatim}}

\keypoint{Why the reference is formulated this way.} The property under test is that
\code{view(B,T,h,d\us{}k).transpose(1,2)} corresponds to contiguous column blocks of the projection.
A reference written with \code{einsum} would be more elegant, but it would share the very
assumption being tested and would therefore prove nothing. Explicit column slicing with a Python
loop is deliberately the least clever formulation available, and that is its merit: it has no
assumptions in common with the implementation.

Observed maximum absolute deviation: $8.94\times10^{-8}$.

\section{Checks 4--5: mathematical invariants}

Checks 1--3 compare against an alternative implementation. Checks 4--5 instead assert properties
that any correct implementation must satisfy, independent of how it is written.

\keypoint{Check 4 --- attention rows sum to one.} Each row of the attention matrix is a probability
distribution over positions. Observed maximum deviation from unity: $1.79\times10^{-7}$. This would
fail immediately if the softmax were applied over the wrong axis --- an error that is otherwise
hard to detect, since the result remains dimensionally valid and the model still trains.

\keypoint{Check 5 --- padded keys receive no attention mass.} The maximum attention weight assigned
to any padded key, over all queries and heads, is $0.00\times10^{0}$ --- exactly zero in floating
point, as the masked logits underflow the softmax completely. This confirms the mask reaches the
scores correctly and is broadcast over the intended axes.

\section{Checks 6--7: padding invariance}
\label{sec:padinv}

These are the decisive checks of the suite, and they test the masking arguments made in
Section~\ref{sec:masking} and in Chapter~\ref{ch:implementation}.

\keypoint{The property.} Appending additional padding tokens to a batch must not change either
model's output. Padding is an artefact of batching and carries no information; if logits move when
padding is added, the model is reading the artefact.

The test constructs a batch with mixed true lengths, appends twenty further padding positions, and
compares logits:

{\small
\begin{verbatim}
short_ids = torch.randint(2, 500, (6, 12))
short_ids[0, 8:] = PAD_ID          # mixed real lengths within the batch
short_ids[3, 5:] = PAD_ID
padded_ids = torch.cat([short_ids, torch.full((6, 20), PAD_ID)], dim=1)
d = (model(short_ids) - model(padded_ids)).abs().max()
\end{verbatim}}

Observed deviations: $0.00\times10^{0}$ for the recurrent model --- exact, as the masked update is
an algebraic identity --- and $1.19\times10^{-7}$ for the Transformer, at the level of
floating-point reassociation.

\keypoint{Why these checks carry the most weight.} The two masking defects described in
Chapter~\ref{ch:implementation} --- unmasked state updates and a backward pass beginning at the
wrong position --- produce no error, no warning and no visible symptom. They degrade accuracy by an
amount that varies with the padding distribution of the batch, which is to say, invisibly. Padding
invariance is the only practical means of detecting them, and the exact-zero result for the
recurrent model confirms that the argument given in Chapter~\ref{ch:implementation} --- that
reverse iteration under masking is equivalent to starting at each sequence's true endpoint --- is
sound in implementation as well as in principle.

\section{Checks 8--9: gradient viability}

The preceding checks all test the forward computation. A model may have a perfectly correct forward
pass and still fail to learn, if gradients do not reach its parameters.

Checks 8 and 9 take a 200-example subset and train each model on it for 30 optimisation steps,
requiring training accuracy above 95\%. A model with a severed gradient path, a detached tensor or
a vanishing initialisation cannot satisfy this.

Both models reach \textbf{100\%} training accuracy within the step budget.

This check also serves a secondary purpose: a model that cannot overfit 200 examples has
insufficient capacity or an optimisation problem, either of which would need to be resolved before
any comparison of generalisation is meaningful.

\section{Results}

\begin{table}[ht]
\centering
\caption[Verification results]{Results of the nine verification checks. All pass. Tolerances are
$10^{-5}$ for equivalence checks, $10^{-6}$ for layer normalisation and $10^{-4}$ for padding
invariance.}
\label{tab:verification}
\begin{tabular}{clll}
\toprule
\textbf{\#} & \textbf{Check} & \textbf{Measured} & \textbf{Result} \\
\midrule
1 & LSTM cell vs.\ reference            & $5.96\times10^{-8}$ & pass \\
2 & Layer norm vs.\ reference           & $4.77\times10^{-7}$ & pass \\
3 & Attention vs.\ naive per-head loop  & $8.94\times10^{-8}$ & pass \\
4 & Attention rows sum to one           & $1.79\times10^{-7}$ & pass \\
5 & Padded keys receive zero mass       & $0.00\times10^{0}$  & pass \\
6 & BiLSTM padding invariance           & $0.00\times10^{0}$  & pass \\
7 & Transformer padding invariance      & $1.19\times10^{-7}$ & pass \\
8 & BiLSTM overfits 200 examples        & 100\% / 30 steps    & pass \\
9 & Transformer overfits 200 examples   & 100\% / 30 steps    & pass \\
\bottomrule
\end{tabular}
\end{table}

All nine checks pass. The implementations are therefore numerically faithful to the mathematics of
Chapter~\ref{ch:theory}, and the experimental results that follow can be attributed to the
architectures rather than to implementation defects.

\section{What this suite does not establish}

Intellectual honesty requires stating the limits of the verification.

The suite establishes that the implemented layers compute the intended functions. It does
\emph{not} establish that the surrounding experimental pipeline is free of error: a defect in the
data splitting, the label mapping or the metric computation would not be detected by any of these
nine checks. Nor does it establish that the chosen architectures are well suited to the task, which
is an empirical question answered in Chapter~\ref{ch:results}.

The balanced-accuracy sanity check --- both models landing in the 90--92\% band predicted by the
literature survey in Table~\ref{tab:prior} --- provides weak additional evidence that the pipeline
as a whole is sound, but it is corroboration rather than proof.
